\section{Sanitizer：这次运行哪里不合法？}

\begin{frame}{Sanitizer是什么？}
    \alert{Sanitizer是一组运行时错误检测工具。}
    \par 编译时为程序加入额外检查；运行时检测到问题，就输出错误类别和位置等诊断信息。
    \medskip
    \begin{columns}[T]
        \begin{column}{0.48\textwidth}
            \begin{block}{ASan：AddressSanitizer}
                检查内存访问错误，例如：
                \begin{itemize}
                    \item 数组越界访问。
                    \item 内存释放后继续使用。
                \end{itemize}
            \end{block}
        \end{column}
        \begin{column}{0.48\textwidth}
            \begin{block}{UBSan：UndefinedBehaviorSanitizer}
                检查特定未定义行为，例如：
                \begin{itemize}
                    \item 有符号整数溢出。
                    \item 整数除以零。
                \end{itemize}
            \end{block}
        \end{column}
    \end{columns}
    \medskip
    \begin{vspquotered}{怎样使用？}
        编译时启用\texttt{-fsanitize=address,undefined}\par
        $\longrightarrow$用测试输入运行程序$\longrightarrow$根据报告定位并修复错误。
    \end{vspquotered}
    \small 配合\texttt{-g}生成调试信息，报告可以显示源码位置；只检查本次实际执行到的操作。
    \note{先建立运行时检测的概念，再进入下一页作业推荐命令。Sanitizer通过编译器插入检查并配合运行库工作，本课使用ASan与UBSan；GDB则用于交互观察状态和执行路径。检测范围与选项有关，没有报告也不代表程序完全正确。}
\end{frame}

\begin{frame}{作业推荐命令里，为什么有Sanitizer？}
    HW1／HW2的README给出了\texttt{-fsanitize=undefined}和\texttt{-fsanitize=address}。
    \begin{itemize}
        \item \alert{ASan}：帮助发现越界访问、释放后使用等内存错误。
        \item \alert{UBSan}：帮助发现有符号溢出等特定未定义行为。
        \item \alert{编译警告}检查可疑写法；运行时检查需要程序实际执行到出错操作。
    \end{itemize}
    \medskip
    \begin{vspquotered}{对照作业要求选择实现}
        HW1 P3不用数组或动态内存；HW2 P3允许系数数组，禁止堆内存和\texttt{math.h}。\par
        本节先用固定系数数组和逐项累加，观察对应的两类报告。
    \end{vspquotered}
\end{frame}

\begin{frame}{HW2 P3：错误演示代码与检查构建}
    \small 课堂调试用错误版：固定$P(x)=1+2x+3x^2$，在$x=2$处求值，正确结果应为\alert{17.000}。
    \begin{columns}[T]
        \begin{column}{0.58\textwidth}
            \lstset{basicstyle=\ttfamily\fontsize{7.4}{8.4}\selectfont,
                numberstyle=\fontsize{4.8}{5.5}\selectfont\color{gray},xleftmargin=6pt}
            \begin{Code}*{c}
                \lstinputlisting[lastline=16]{code/sanitizer/polynomial.c}
            \end{Code}
            \scriptsize 完整源码：\texttt{code/sanitizer/polynomial.c}。
        \end{column}
        \begin{column}{0.38\textwidth}
            \small 在\texttt{r4/}中编译：
            \lstset{basicstyle=\ttfamily\fontsize{6.5}{7.8}\selectfont,
                numberstyle=\fontsize{4.8}{5.5}\selectfont\color{gray},xleftmargin=6pt,
                breakatwhitespace=true}
            \begin{Code}*{bash}
                \lstinputlisting{code/commands/build-sanitizer.sh}
            \end{Code}
            \scriptsize
            \begin{itemize}
                \item \texttt{-g -O0}便于定位源码与逐行观察。
                \item \texttt{-fsanitize=address,undefined}启用ASan与UBSan，链接时也需要。
                \item UBSan报告错误后默认通常会继续执行；ASan发现此类内存错误后通常会终止程序。
            \end{itemize}
        \end{column}
    \end{columns}
    \note{完整引用提交包中的code/sanitizer/polynomial.c。保留i=count的错误，首次访问coef[3]，由ASan报告栈数组越界；改为count-1后输出17.000。这是依据HW2 P3要求编写的固定数据课堂变体，不假定完整题面的输入格式；正式评测参数以作业说明为准。}
\end{frame}

\begin{frame}{HW2 P3：解读这份ASan实测报告}
    \begin{columns}[T]
        \begin{column}{0.61\textwidth}
            \small 终端实测报告节选（路径统一为包内路径）：
            \lstset{basicstyle=\ttfamily\fontsize{6.4}{8.0}\selectfont,
                numbers=none,xleftmargin=0pt,breakatwhitespace=true}
            \begin{Code}*{plain}
                \lstinputlisting[linerange={2-5,9-13}]{code/output/asan-hw2-p3-user.txt}
            \end{Code}
            \scriptsize 完整报告保留于课件源码；地址与进程号随运行变化。
        \end{column}
        \begin{column}{0.35\textwidth}
            \footnotesize
            \begin{enumerate}
                \item \alert{错误类别}\par
                    \texttt{stack-buffer-overflow}：栈上数组越界。
                \item \alert{读取与位置}\par
                    读取8字节，即本例一个\texttt{double}；出错于第7行\texttt{coef[i]}。
                \item \alert{调用栈}\par
                    \texttt{main}第14行调用\texttt{evaluate}，在其中触发错误。
                \item \alert{数组范围}\par
                    \texttt{[32, 56)}共24字节；偏移56恰好越过末尾，对应\texttt{coef[3]}。
            \end{enumerate}
            \medskip
            \alert{循环从count - 1开始。}\par
            重新编译并运行，预期输出\texttt{17.000}。
        \end{column}
    \end{columns}
    \note{报告选自用户实测，保存为code/output/asan-hw2-p3-user.txt。为课堂直接定位源码，文本中的路径统一为提交包内相对路径，地址和诊断内容保留实测值；页面选取第2—5与9—13行。先看错误类别、访问大小、源码位置、调用栈，再看对象范围；[32,56)是相对于main栈帧的偏移区间，三个double共24字节。Shadow bytes是补充诊断，本次先聚焦数组下标和循环初值。}
\end{frame}

\begin{frame}{HW2 P3：多项式求值为什么越界了？}
    \begin{columns}[T]
        \begin{column}{0.61\textwidth}
            \begin{Code}*{c}
                \lstinputlisting[firstline=4,lastline=10,firstnumber=4]{code/sanitizer/polynomial.c}
            \end{Code}
        \end{column}
        \begin{column}{0.35\textwidth}
            课堂固定数据：\par
            \texttt{coef = \{1, 2, 3\}}\par
            $P(x)=1+2x+3x^2$\par
            \texttt{count = 3}\par
            \texttt{x = 2}\par\medskip
            合法结果为\alert{17.000}。
        \end{column}
    \end{columns}
    \medskip
    从高次项向低次项递推：\texttt{value = value * x + coef[i]}。\par
    \alert{第一次循环中，i的值是否是合法下标？}
    \note{课堂变体仅固定一个多项式，不假定正式题面的输入顺序或完整范围。本节约22分钟，含7至8分钟演示。}
\end{frame}

\begin{frame}{三个系数，合法下标只有0、1、2}
    \begin{center}
        \begin{tikzpicture}[>=Stealth,cell/.style={draw=VSPPrimary,
            minimum width=1.8cm,minimum height=1.05cm,font=\large}]
            \node[cell] (a) at (0,0) {1};
            \node[cell] (b) at (1.8,0) {2};
            \node[cell] (c) at (3.6,0) {3};
            \node[cell,dashed,fill=red!12] (d) at (5.4,0) {?};
            \node[below=0.15cm of a] {\texttt{coef[0]}};
            \node[below=0.15cm of b] {\texttt{coef[1]}};
            \node[below=0.15cm of c] {\texttt{coef[2]}};
            \node[below=0.15cm of d,text=VSPPrimary] {\texttt{coef[3]}};
            \draw[->,thick,VSPPrimary] (5.4,1.25)--(d.north)
                node[pos=0,above,font=\small] {第一次读取就越界};
        \end{tikzpicture}
    \end{center}
    \medskip
    \begin{itemize}
        \item 初值\texttt{i = count}使第一次访问成为\texttt{coef[3]}。
        \item 修复初值为\texttt{count - 1}，依次访问下标2、1、0。
        \item 本课堂变体使用栈上的固定数组；\alert{没有引入堆分配或VLA}。
    \end{itemize}
\end{frame}

\begin{frame}{第一次读懂ASan报告}
    \begin{Code}*{plain}
        \lstinputlisting{code/output/asan-bounds.txt}
    \end{Code}
    \begin{enumerate}
        \item \alert{错误类别}：\texttt{stack-buffer-overflow}，栈上的越界访问。
        \item \alert{源码位置}：在\texttt{evaluate}函数读取\texttt{coef[i]}时出错。
        \item \alert{对象范围}：回到\texttt{main}中的固定数组，核对长度与循环初值。
    \end{enumerate}
    \small 此处为本课程序实测报告节选；省略地址、进程号和部分调用栈。
\end{frame}

\begin{frame}{power：乘法在哪一轮超出范围？}
    \begin{columns}[T]
        \begin{column}{0.48\textwidth}
            \small 固定计算\texttt{power(50000, 2)}：
            \lstset{basicstyle=\ttfamily\fontsize{7.4}{8.7}\selectfont,
                numberstyle=\tiny\color{gray},xleftmargin=6pt}
            \begin{Code}*{c}
                \lstinputlisting[lastline=14]{code/sanitizer/power.c}
            \end{Code}
            \small 第一次乘法得到50000；第二次的数学结果为\alert{2500000000}。
        \end{column}
        \begin{column}{0.48\textwidth}
            \small 在\texttt{slides/r4/}中启用UBSan：
            \lstset{basicstyle=\ttfamily\fontsize{6.4}{7.8}\selectfont,
                numbers=none,xleftmargin=0pt,breakatwhitespace=true}
            \begin{Code}*{bash}
                \lstinputlisting{code/commands/power-sanitizer.sh}
            \end{Code}
            \small 实测诊断（第6行）：
            \begin{Code}*{plain}
                \lstinputlisting{code/output/ubsan-power.txt}
            \end{Code}
            \footnotesize 当前32位\texttt{int}上界是2147483647，第二次乘法触发有符号溢出UB。
        \end{column}
    \end{columns}
    \note{本例使用固定参数，先让学生手算两轮结果。UBSan在result *= x处报告50000 * 50000不能由int表示。原程序与修复版演示副本位于工作区demo/power/，不要依赖原程序关闭检测后出现的某个回绕值。}
\end{frame}

\begin{frame}{power：GDB定位与修复验证}
    \begin{columns}[T]
        \begin{column}{0.46\textwidth}
            \small 启动\texttt{gdb -q build/demo/power-ubsan}，在乘法执行前停住：
            \lstset{basicstyle=\ttfamily\fontsize{7.2}{8.5}\selectfont,
                numbers=none,xleftmargin=0pt,breakatwhitespace=true}
            \begin{Code}*{plain}
                \lstinputlisting{code/commands/power-session.gdb}
            \end{Code}
            \small 实测\texttt{i=1}、\texttt{result=50000}、\texttt{x=50000}。
            \par\medskip
            \texttt{bt}显示当前在\texttt{power}，由\texttt{main}调用；下一次乘法超出\texttt{int}范围。
        \end{column}
        \begin{column}{0.50\textwidth}
            \small 修复乘积变量、返回类型与输出格式：
            \lstset{basicstyle=\ttfamily\fontsize{7.2}{8.5}\selectfont,
                numbers=left,numberstyle=\tiny\color{gray},xleftmargin=6pt}
            \begin{Code}*{c}
                \lstinputlisting[lastline=14]{code/fixed/power.c}
            \end{Code}
            \footnotesize 使用相同检查选项重新编译，输出\alert{2500000000}，无UBSan诊断。
        \end{column}
    \end{columns}
    \note{条件断点位于power.c第6行，i==1表示第二次循环，断点停在执行乘法之前。本例结果可以由long long表示；宽类型并不保证任意底数与指数都不溢出，仍应根据输入及中间结果范围设计实现。完整修复版为code/fixed/power.c与工作区demo/power/power_fixed.c。GDB批处理脚本与变量值已经实测验证。}
\end{frame}

\begin{frame}{数组循环：最后一次迭代访问了哪里？}
    \begin{columns}[T]
        \begin{column}{0.48\textwidth}
            \small 目标：让6个数组元素分别加1。
            \lstset{basicstyle=\ttfamily\fontsize{7.5}{9.0}\selectfont,
                numberstyle=\tiny\color{gray},xleftmargin=6pt}
            \begin{Code}*{c}
                \lstinputlisting[lastline=10]{code/sanitizer/array_loop.c}
            \end{Code}
            \small 循环包含\texttt{i == 6}；此时访问\texttt{data[6]}，越过数组末尾。
        \end{column}
        \begin{column}{0.48\textwidth}
            \small 在\texttt{slides/r4/}中单独启用ASan：
            \lstset{basicstyle=\ttfamily\fontsize{6.4}{7.8}\selectfont,
                numbers=none,xleftmargin=0pt,breakatwhitespace=true}
            \begin{Code}*{bash}
                \lstinputlisting{code/commands/array-sanitizer.sh}
            \end{Code}
            \small 实测报告节选：
            \begin{Code}*{plain}
                \lstinputlisting{code/output/asan-array-loop.txt}
            \end{Code}
            \footnotesize 自增先读后写，本次先报告读取4字节；\texttt{data}范围为\texttt{[32, 56)}，偏移56已越界。
        \end{column}
    \end{columns}
    \note{这是固定长度的栈数组例子。报告来自本程序实际执行，完整输出位于build/demo/output/array-loop-asan.txt，页面省略部分地址与调用栈。本页单独使用ASan以观察内存访问报告；同时启用UBSan时，也可能先出现数组下标越界诊断。工作区demo/array/提供同内容源码副本。}
\end{frame}

\begin{frame}{数组循环：用条件断点验证修复}
    \begin{columns}[T]
        \begin{column}{0.46\textwidth}
            \small 启动\texttt{gdb -q build/demo/array-loop-asan}，在越界访问之前停止：
            \lstset{basicstyle=\ttfamily\fontsize{7.2}{8.5}\selectfont,
                numbers=none,xleftmargin=0pt,breakatwhitespace=true}
            \begin{Code}*{plain}
                \lstinputlisting{code/commands/array-session.gdb}
            \end{Code}
            \small 实测\texttt{i=6}，数组内容为6个1，\texttt{sizeof data=24}。
            \par\medskip
            前6轮已完成全部元素的更新；\alert{第7次迭代多执行了一次。}
        \end{column}
        \begin{column}{0.50\textwidth}
            \small 修复：将循环条件改为\texttt{i < 6}。
            \lstset{basicstyle=\ttfamily\fontsize{7.5}{9.0}\selectfont,
                numbers=left,numberstyle=\tiny\color{gray},xleftmargin=6pt}
            \begin{Code}*{c}
                \lstinputlisting[lastline=10]{code/fixed/array_loop.c}
            \end{Code}
            \small 重新编译运行，输出\alert{1}，无ASan诊断。
            \par\medskip
            \footnotesize 完整修复版：\par\texttt{code/fixed/array\_loop.c}。
        \end{column}
    \end{columns}
    \note{条件断点在array_loop.c第6行、i==6时停住，尚未访问越界元素。set print repeats 0让GDB展开显示重复数组值；当前数组六个元素均为1，说明正常更新已经完成，循环边界导致多执行一步。修复后验证退出码0、输出1且无Sanitizer报告；GDB脚本已实测。}
\end{frame}


\begin{frame}{没有报告，还需要检查什么？}
    \begin{itemize}
        \item \alert{测试覆盖}：本次没执行到的代码，运行时检查没有替你验证。
        \item \alert{工具覆盖}：ASan和UBSan各有检测范围，不能覆盖所有语言错误。
        \item \alert{业务逻辑}：最大值更新用了总和、漏算常数项，都可能没有报告。
        \item \alert{修复验证}：重新编译，重跑失败输入，再加正常与边界输入。
    \end{itemize}
    \medskip
    \begin{vspquotered}{把工具放进同一个调试过程}
        Sanitizer帮助定位非法操作，GDB帮助观察状态与执行路径。\par
        结合语言规则与题目要求，解释为什么出错、为什么修复有效。
    \end{vspquotered}
\end{frame}
